\documentclass[UTF8,zihao=-4]{ctexart}
\usepackage[a4paper,margin=2.2cm]{geometry}
\usepackage{hyperref}
\usepackage{xcolor}
\usepackage{enumitem}
\usepackage{fancyvrb}
\usepackage{fvextra}
\usepackage{longtable}
\hypersetup{colorlinks=true,linkcolor=blue,urlcolor=blue}
\setlist{nosep,leftmargin=2em}
\setlength{\parskip}{0.45em}
\setlength{\parindent}{0pt}
\title{企业知识库 Agent 助手项目完整帮助文档}
\author{基于当前项目代码实时更新}
\date{2026-05-09}

\begin{document}
\maketitle
\tableofcontents
\newpage

\section{项目一句话说明}

这是一个基于 \texttt{Streamlit + LangChain ReAct Agent + Chroma + DashScope} 的企业内部知识库智能助手示例项目。项目面向 AI 应用开发学习、简历展示和企业知识问答演示，支持企业制度问答、员工模拟数据查询、个人周报生成、知识文档上传入库和流式对话展示。

当前项目已经从早期的产品客服场景迁移为企业知识库场景。知识库资料位于 \texttt{data/enterprise/}，结构化员工模拟数据位于 \texttt{data/enterprise\_external/employee\_records.csv}，主界面入口是 \texttt{app.py}。

\section{项目能做什么}

当前项目主要支持以下能力：

\begin{enumerate}
\item 企业制度知识库问答：基于 HR、财务、IT 支持、项目协作、信息安全、入职指引等制度文档进行 RAG 检索增强回答。
\item ReAct Agent 工具调度：Agent 根据用户问题自动选择知识库检索、员工 ID 获取、部门查询、员工记录查询、周报上下文生成等工具。
\item 员工模拟数据查询：支持查询年假余额、报销状态、项目记录、入职待办和周报相关业务数据。
\item 个人周报生成：按固定工具流程获取员工 ID、部门、周报上下文和项目记录，再生成结构化中文周报。
\item Streamlit 对话界面：提供类聊天体验、流式输出、检索来源折叠展示、示例问题入口、模型和检索参数控制。
\item 知识库维护：侧边栏可上传 TXT/PDF 文档，保存到 \texttt{data/enterprise/} 并写入 Chroma 向量库。
\item 会话导出：可将当前对话导出为 Markdown 文件。
\end{enumerate}

\section{真实调用链路}

\begin{Verbatim}[breaklines=true,breakanywhere=true,fontsize=\small]
Streamlit(app.py)
  -> ReactAgent(agent/react_agent.py)
  -> AgentExecutor.stream()
  -> Tools(agent/tools/agent_tools.py)
  -> RagSummarizeService(rag/rag_service.py)
  -> VectorStoreService(rag/vector_store.py)
  -> Chroma(chroma_db/)
  -> DashScope(model/factory.py)
  -> Prompt(prompts/*.txt)
\end{Verbatim}

用户在 \texttt{app.py} 的 \texttt{st.chat\_input()} 输入问题后，前端会调用 \texttt{ReactAgent.execute\_stream()}。Agent 内部通过 \texttt{create\_react\_agent()} 创建 ReAct Agent，再由 \texttt{AgentExecutor} 执行工具调用循环。最终答案通过 \texttt{AgentExecutor.stream()} 分块返回给 Streamlit，并由 \texttt{st.write\_stream()} 展示。

\section{目录结构}

\begin{Verbatim}[breaklines=true,breakanywhere=true,fontsize=\small]
.
|-- app.py                         # Streamlit 应用入口
|-- requirements.txt               # Python 运行依赖
|-- help_requirements.txt          # 帮助文档生成相关依赖
|-- HELP_PROJECT_GUIDE.tex         # 当前帮助文档 LaTeX 源文件
|-- HELP_PROJECT_GUIDE.pdf         # 当前帮助文档 PDF
|-- build_help_pdf.py              # 旧版 Markdown 到 LaTeX 辅助脚本
|-- build_help_pdf_reportlab.py    # 旧版 ReportLab PDF 辅助脚本
|-- md5_enterprise.txt             # 已入库知识文件 MD5 记录
|-- config/
|   |-- agent.yml                  # 员工模拟数据路径配置
|   |-- chroma.yml                 # Chroma 与知识库入库配置
|   |-- prompts.yml                # Prompt 文件路径配置
|   `-- rag.yml                    # DashScope 模型配置
|-- agent/
|   |-- react_agent.py             # ReAct Agent 创建与流式执行逻辑
|   `-- tools/
|       |-- agent_tools.py         # 企业场景工具函数
|       `-- middleware.py          # 中间件示例代码
|-- rag/
|   |-- vector_store.py            # 文档加载、切分、向量入库、Retriever 创建
|   `-- rag_service.py             # RAG 检索与总结链路
|-- model/
|   `-- factory.py                 # ChatTongyi 与 DashScope Embeddings 初始化
|-- prompts/
|   |-- main_prompt.txt            # Agent 主 Prompt
|   |-- rag_summarize.txt          # RAG 总结 Prompt
|   `-- report_prompt.txt          # 报告生成 Prompt
|-- data/
|   |-- enterprise/                # 企业制度知识库文档
|   |   |-- finance_reimbursement_policy.txt
|   |   |-- hr_leave_policy.txt
|   |   |-- it_support_policy.txt
|   |   |-- onboarding_guide.txt
|   |   |-- project_collaboration_policy.txt
|   |   `-- security_policy.txt
|   `-- enterprise_external/
|       `-- employee_records.csv   # 员工模拟结构化数据
|-- utils/                         # 配置、路径、文件、日志工具
|-- chroma_db/                     # Chroma 本地持久化目录
`-- logs/                          # 运行日志目录
\end{Verbatim}

\section{环境要求与安装}

运行项目前需要准备：

\begin{itemize}
\item Python 3.10 或更高版本。
\item 可用的阿里云 DashScope API Key。
\item Windows PowerShell、CMD、macOS 或 Linux Shell 均可运行。
\end{itemize}

安装依赖：

\begin{Verbatim}[breaklines=true,breakanywhere=true,fontsize=\small]
pip install -r requirements.txt
\end{Verbatim}

如果下载较慢，可使用清华源：

\begin{Verbatim}[breaklines=true,breakanywhere=true,fontsize=\small]
pip install -r requirements.txt -i https://pypi.tuna.tsinghua.edu.cn/simple
\end{Verbatim}

当前 \texttt{requirements.txt} 中的关键依赖包括：

\begin{Verbatim}[breaklines=true,breakanywhere=true,fontsize=\small]
langchain==0.3.7
langchain-core==0.3.18
langchain-community==0.3.7
langchain-chroma==0.1.4
chromadb==0.5.15
langchain-text-splitters==0.3.2
langgraph==0.2.50
streamlit==1.40.1
pypdf==5.1.0
pyyaml==6.0.2
dashscope==1.20.14
\end{Verbatim}

\section{配置 DashScope API Key}

项目通过 DashScope 调用通义千问聊天模型和 Embedding 模型，运行前需要配置 \texttt{DASHSCOPE\_API\_KEY}。

PowerShell 临时配置：

\begin{Verbatim}[breaklines=true,breakanywhere=true,fontsize=\small]
$env:DASHSCOPE_API_KEY="your-api-key"
\end{Verbatim}

CMD 临时配置：

\begin{Verbatim}[breaklines=true,breakanywhere=true,fontsize=\small]
set DASHSCOPE_API_KEY=your-api-key
\end{Verbatim}

macOS/Linux 临时配置：

\begin{Verbatim}[breaklines=true,breakanywhere=true,fontsize=\small]
export DASHSCOPE_API_KEY="your-api-key"
\end{Verbatim}

\section{核心配置说明}

\subsection{模型配置}

文件：\texttt{config/rag.yml}

\begin{Verbatim}[breaklines=true,breakanywhere=true,fontsize=\small]
chat_model_name: qwen3-max
embedding_model_name: text-embedding-v4
\end{Verbatim}

\begin{itemize}
\item \texttt{chat\_model\_name}：聊天模型名称，当前使用 \texttt{qwen3-max}。
\item \texttt{embedding\_model\_name}：向量化模型名称，当前使用 \texttt{text-embedding-v4}。
\end{itemize}

\subsection{向量库配置}

文件：\texttt{config/chroma.yml}

\begin{Verbatim}[breaklines=true,breakanywhere=true,fontsize=\small]
collection_name : enterprise_agent
persist_directory : chroma_db
k : 3
data_path : data/enterprise
md5_hex_store : md5_enterprise.txt
allow_knowledge_file_type : ["txt","pdf"]
chunk_size : 200
chunk_overlap : 20
\end{Verbatim}

\begin{itemize}
\item \texttt{collection\_name}：Chroma 集合名称，当前为 \texttt{enterprise\_agent}。
\item \texttt{persist\_directory}：向量库持久化目录，当前为 \texttt{chroma\_db}。
\item \texttt{k}：每次检索返回的 top-k 文档片段数，当前为 3。
\item \texttt{data\_path}：知识库文档目录，当前为 \texttt{data/enterprise}。
\item \texttt{md5\_hex\_store}：已处理文件 MD5 记录文件，当前为 \texttt{md5\_enterprise.txt}。
\item \texttt{allow\_knowledge\_file\_type}：允许入库的文档类型，当前支持 \texttt{txt} 和 \texttt{pdf}。
\item \texttt{chunk\_size} 和 \texttt{chunk\_overlap}：控制文档切分粒度和相邻片段重叠长度。
\end{itemize}

\subsection{Prompt 配置}

文件：\texttt{config/prompts.yml}

\begin{Verbatim}[breaklines=true,breakanywhere=true,fontsize=\small]
main_prompt_path : prompts/main_prompt.txt
rag_summarize_prompt_path : prompts/rag_summarize.txt
report_prompt_path : prompts/report_prompt.txt
\end{Verbatim}

\begin{itemize}
\item \texttt{main\_prompt.txt}：Agent 主提示词，定义企业知识库助手角色、工具边界和固定工作流。
\item \texttt{rag\_summarize.txt}：RAG 检索后的总结提示词，要求答案基于参考资料。
\item \texttt{report\_prompt.txt}：报告生成提示词，目前保留为报告场景扩展入口。
\end{itemize}

\subsection{员工模拟数据配置}

文件：\texttt{config/agent.yml}

\begin{Verbatim}[breaklines=true,breakanywhere=true,fontsize=\small]
external_data_path: data/enterprise_external/employee_records.csv
\end{Verbatim}

\texttt{employee\_records.csv} 模拟企业业务系统中的员工数据，字段包括：

\begin{Verbatim}[breaklines=true,breakanywhere=true,fontsize=\small]
employee_id, employee_name, department, leave_balance_days,
reimbursement_status, missing_reimbursement_docs, project_name,
weekly_highlights, deliverables, next_week_plan, risks, onboarding_tasks
\end{Verbatim}

\section{核心知识点}

\subsection{什么是 RAG}

RAG 是 Retrieval-Augmented Generation，即检索增强生成。它先从外部知识库检索和问题相关的资料，再把资料连同用户问题一起交给大模型生成答案。

本项目中的 RAG 链路是：

\begin{Verbatim}[breaklines=true,breakanywhere=true,fontsize=\small]
用户问题 -> Chroma 检索企业制度片段 -> 拼接 context -> DashScope 模型总结回答
\end{Verbatim}

它解决的问题是：大模型自身不知道本地企业制度文档、流程规范和模拟业务资料，所以需要先检索，再回答。

\subsection{为什么要切分文档}

原始制度文档可能较长，直接整体向量化会降低检索粒度。切分为 chunk 后，每个片段都可以单独向量化和检索，更容易命中具体条款、流程步骤或注意事项。

\subsection{什么是 Chroma}

Chroma 是本地向量数据库，用来存储文档片段、向量和元数据。项目使用 \texttt{langchain\_chroma.Chroma} 创建本地持久化向量库：

\begin{Verbatim}[breaklines=true,breakanywhere=true,fontsize=\small]
Chroma(
    collection_name="enterprise_agent",
    embedding_function=embed_model,
    persist_directory="chroma_db",
)
\end{Verbatim}

\subsection{什么是 ReAct Agent}

ReAct 是 Reasoning + Acting。模型按照“思考、选择工具、读取观察结果、继续推理、最终回答”的循环工作。普通 LLM 调用只生成文本，ReAct Agent 可以根据问题选择工具并把工具结果整合成最终答案。

本项目中的 Agent 使用以下标准格式：

\begin{Verbatim}[breaklines=true,breakanywhere=true,fontsize=\small]
Thought -> Action -> Action Input -> Observation -> Final Answer
\end{Verbatim}

\section{Agent 工具说明}

当前 \texttt{agent/react\_agent.py} 注册的工具列表为：

\begin{Verbatim}[breaklines=true,breakanywhere=true,fontsize=\small]
tools = [
    rag_summarize,
    get_employee_id,
    get_employee_department,
    fetch_employee_records,
    generate_weekly_report_context,
]
\end{Verbatim}

\begin{longtable}{p{0.28\textwidth}p{0.62\textwidth}}
\textbf{工具} & \textbf{用途} \\
\hline
\texttt{rag\_summarize} & 从企业知识库检索制度、流程、IT 支持、项目规范、信息安全等参考资料，并生成基于资料的回答。 \\
\texttt{get\_employee\_id} & 获取当前员工 ID。当前为演示逻辑，会从 \texttt{E1001} 到 \texttt{E1005} 中随机返回一个。 \\
\texttt{get\_employee\_department} & 根据员工 ID 查询员工所属部门。 \\
\texttt{fetch\_employee\_records} & 查询员工模拟业务数据，入参格式为 \texttt{员工ID,查询类型}。查询类型包括 \texttt{leave\_balance}、\texttt{reimbursement}、\texttt{projects}、\texttt{onboarding}、\texttt{weekly\_report}、\texttt{all}。 \\
\texttt{generate\_weekly\_report\_context} & 进入员工周报生成场景，提示 Agent 输出本周概览、关键产出、问题风险和下周计划。 \\
\end{longtable}

\section{主要工作流程}

\subsection{企业制度问答}

\begin{Verbatim}[breaklines=true,breakanywhere=true,fontsize=\small]
用户询问制度、流程或规范
  -> Streamlit 调用 ReactAgent.execute_stream()
  -> Agent 判断需要企业知识库资料
  -> 调用 rag_summarize
  -> RagSummarizeService 检索 Chroma
  -> 拼接参考资料 context
  -> DashScope 模型基于资料生成回答
  -> Streamlit 流式展示答案并折叠展示检索来源
\end{Verbatim}

\subsection{员工个人数据查询}

当用户询问“我的年假还有多少”“我的报销状态如何”“我有哪些入职待办”等个人数据问题时，Agent 应按以下流程：

\begin{enumerate}
\item 如果用户没有显式提供员工 ID，先调用 \texttt{get\_employee\_id}。
\item 调用 \texttt{fetch\_employee\_records}，入参为 \texttt{员工ID,查询类型}。
\item 只根据工具返回的模拟结构化数据回答，不编造额外状态。
\end{enumerate}

常见查询类型：

\begin{Verbatim}[breaklines=true,breakanywhere=true,fontsize=\small]
leave_balance   # 年假余额
reimbursement   # 报销状态
projects        # 项目记录
onboarding      # 入职待办
weekly_report   # 周报数据
all             # 返回完整记录
\end{Verbatim}

\subsection{个人周报生成}

周报生成在 \texttt{agent/react\_agent.py} 中被写入固定流程：

\begin{enumerate}
\item 调用 \texttt{get\_employee\_id} 获取员工 ID。
\item 调用 \texttt{get\_employee\_department} 获取部门。
\item 调用 \texttt{generate\_weekly\_report\_context} 进入周报生成场景。
\item 调用 \texttt{fetch\_employee\_records}，入参为 \texttt{employee\_id,weekly\_report}。
\item 用中文生成结构化周报。
\end{enumerate}

\subsection{知识文档上传和入库}

Streamlit 侧边栏提供 TXT/PDF 上传入口。点击“保存并载入知识库”后，文件会保存到 \texttt{data/enterprise/}，然后调用：

\begin{Verbatim}[breaklines=true,breakanywhere=true,fontsize=\small]
rag_service.vector_store.load_document()
\end{Verbatim}

入库过程会扫描 \texttt{data/enterprise/}，按 \texttt{config/chroma.yml} 的配置读取 TXT/PDF，计算 MD5 去重，切分文档，生成 Embedding，并写入 \texttt{chroma\_db/}。

\section{启动应用}

在项目根目录执行：

\begin{Verbatim}[breaklines=true,breakanywhere=true,fontsize=\small]
streamlit run app.py
\end{Verbatim}

启动成功后，终端通常会显示：

\begin{Verbatim}[breaklines=true,breakanywhere=true,fontsize=\small]
Local URL: http://localhost:8501
\end{Verbatim}

浏览器访问该地址即可使用企业知识库问答界面。

\section{推荐测试问题}

企业制度问答：

\begin{Verbatim}[breaklines=true,breakanywhere=true,fontsize=\small]
差旅报销需要哪些材料？
公司对外部 AI 工具使用有什么安全要求？
VPN 无法连接应该怎么处理？
新员工入职第一周应该完成哪些事项？
\end{Verbatim}

员工数据查询：

\begin{Verbatim}[breaklines=true,breakanywhere=true,fontsize=\small]
我还剩多少年假？
我的报销状态怎么样？
我有哪些入职待办？
我本周参与了哪些项目工作？
\end{Verbatim}

周报生成：

\begin{Verbatim}[breaklines=true,breakanywhere=true,fontsize=\small]
根据我的本周项目记录生成周报。
帮我整理一份个人本周工作总结。
\end{Verbatim}

\section{主要代码入口}

\subsection{\texttt{app.py}}

\texttt{app.py} 是 Streamlit 前端入口，主要职责包括：

\begin{itemize}
\item 初始化 \texttt{ReactAgent} 和会话状态。
\item 注入自定义 CSS，构建企业知识库对话界面。
\item 在侧边栏提供模型选择、Temperature、Top-K 检索数量、文件上传、会话导出和清空会话。
\item 接收用户输入，调用 \texttt{stream\_agent\_response()}。
\item 预先检索来源文档，并在助手回复后以折叠面板展示引用片段。
\item 将完整助手回复通过 \texttt{"".join(response\_chunks)} 保存到会话历史。
\end{itemize}

\subsection{\texttt{agent/react\_agent.py}}

负责构建 ReAct Agent：

\begin{itemize}
\item 读取 \texttt{prompts/main\_prompt.txt}。
\item 注册企业场景工具列表。
\item 拼接 ReAct 标准模板和周报、员工数据固定流程。
\item 使用 \texttt{create\_react\_agent()} 创建 Agent。
\item 使用 \texttt{AgentExecutor} 包装执行，并启用 \texttt{handle\_parsing\_errors}。
\item 通过 \texttt{execute\_stream()} 提供流式输出。
\end{itemize}

\subsection{\texttt{agent/tools/agent\_tools.py}}

定义 Agent 可调用工具。该文件会初始化 \texttt{RagSummarizeService}，并通过 \texttt{@tool} 暴露企业知识检索、员工 ID、部门、员工记录和周报上下文工具。

\subsection{\texttt{rag/vector\_store.py}}

负责知识库入库和 retriever 创建：

\begin{itemize}
\item 初始化 \texttt{Chroma}。
\item 使用 \texttt{RecursiveCharacterTextSplitter} 切分文档。
\item 扫描 \texttt{data/enterprise/} 中允许类型的文档。
\item 用 MD5 记录避免重复入库。
\item 通过 \texttt{add\_documents()} 写入向量库。
\item 返回 \texttt{as\_retriever(search\_kwargs=\{"k": 3\})}。
\end{itemize}

\subsection{\texttt{rag/rag\_service.py}}

负责 RAG 问答链路：

\begin{itemize}
\item 调用 retriever 检索相关文档。
\item 拼接参考资料 \texttt{context}。
\item 使用 \texttt{rag\_summarize.txt} 限制模型基于资料回答。
\item 通过 \texttt{StrOutputParser()} 输出纯字符串。
\end{itemize}

\subsection{\texttt{model/factory.py}}

负责模型初始化：

\begin{itemize}
\item \texttt{ChatTongyi(model=rag\_conf["chat\_model\_name"])}：聊天模型。
\item \texttt{DashScopeEmbeddings(model=rag\_conf["embedding\_model\_name"])}：Embedding 模型。
\end{itemize}

\section{如何新增知识库文档}

\begin{enumerate}
\item 将 \texttt{.txt} 或 \texttt{.pdf} 文件放入 \texttt{data/enterprise/}。
\item 确认 \texttt{config/chroma.yml} 中允许该文件类型。
\item 执行知识库入库命令：
\end{enumerate}

\begin{Verbatim}[breaklines=true,breakanywhere=true,fontsize=\small]
python -m rag.vector_store
\end{Verbatim}

也可以在 Streamlit 侧边栏上传 TXT/PDF 后点击“保存并载入知识库”。

建议：

\begin{itemize}
\item TXT 文件使用 UTF-8 编码。
\item 文档内容按主题分段，便于切分和检索。
\item 如果问题经常检索不到，可适当调大 Top-K，或优化知识库文档表达。
\item 如果文档改动后没有重新入库，检查 \texttt{md5\_enterprise.txt} 和 \texttt{logs/}。
\item 如果确实需要完全重建向量库，请手动确认旧向量库和 MD5 记录的处理方式，不要使用批量递归删除命令。
\end{itemize}

\section{如何新增 Agent 工具}

新增工具通常需要修改两处：\texttt{agent/tools/agent\_tools.py} 和 \texttt{agent/react\_agent.py}。

第一步，在 \texttt{agent/tools/agent\_tools.py} 中定义工具：

\begin{Verbatim}[breaklines=true,breakanywhere=true,fontsize=\small]
from langchain_core.tools import tool

@tool
def get_policy_owner(policy_name: str) -> str:
    """根据制度名称查询制度负责人。入参为制度名称字符串。"""
    return "HR"
\end{Verbatim}

第二步，在 \texttt{agent/react\_agent.py} 中导入并加入 \texttt{tools} 列表：

\begin{Verbatim}[breaklines=true,breakanywhere=true,fontsize=\small]
tools = [
    rag_summarize,
    get_employee_id,
    get_policy_owner,
]
\end{Verbatim}

第三步，在 \texttt{prompts/main\_prompt.txt} 中补充工具调用场景和入参格式。工具 docstring 和 Prompt 描述越清楚，Agent 越容易正确调用。

\section{常见问题}

\subsection{启动时报 DashScope API Key 错误}

检查是否设置了环境变量：

\begin{Verbatim}[breaklines=true,breakanywhere=true,fontsize=\small]
$env:DASHSCOPE_API_KEY
\end{Verbatim}

如果没有输出，重新设置：

\begin{Verbatim}[breaklines=true,breakanywhere=true,fontsize=\small]
$env:DASHSCOPE_API_KEY="your-api-key"
\end{Verbatim}

\subsection{问答结果和知识库内容无关}

建议检查：

\begin{itemize}
\item \texttt{chroma\_db/} 是否已经生成。
\item 是否执行过 \texttt{python -m rag.vector\_store}。
\item \texttt{data/enterprise/} 中是否有相关制度文档。
\item \texttt{config/chroma.yml} 中的 \texttt{k} 是否过小。
\item 文档编码是否为 UTF-8。
\item \texttt{rag\_summarize.txt} 是否仍要求模型基于参考资料回答。
\end{itemize}

\subsection{新增文档后没有被检索到}

项目使用 MD5 记录已处理文件。如果文件内容没有变化，再次运行入库时会跳过相同 MD5 的文件。可以查看 \texttt{md5\_enterprise.txt} 和 \texttt{logs/} 确认是否被跳过。

\subsection{Agent 输出格式错误或中断}

ReAct Agent 对输出格式较严格。项目已在 \texttt{AgentExecutor} 中配置 \texttt{handle\_parsing\_errors}。如果仍失败，可以：

\begin{itemize}
\item 换一种更明确的问题问法。
\item 检查 \texttt{prompts/main\_prompt.txt} 是否被改坏。
\item 换用更稳定的聊天模型。
\item 减少 Prompt 中和 ReAct 标准格式冲突的说明。
\end{itemize}

\subsection{中文显示乱码}

如果 Windows 终端中显示乱码，优先用支持 UTF-8 的编辑器查看文件。LaTeX 源文件应保存为 UTF-8，并使用支持中文的编译方式生成 PDF。

\section{学习路线}

建议按以下顺序阅读代码：

\begin{enumerate}
\item 先看 \texttt{app.py}，理解前端如何接收问题、展示答案、上传文档和导出会话。
\item 再看 \texttt{agent/react\_agent.py}，理解 Agent 如何创建、如何注册工具、如何约束周报流程。
\item 再看 \texttt{agent/tools/agent\_tools.py}，理解每个工具具体做什么。
\item 再看 \texttt{rag/vector\_store.py}，理解文档如何入库。
\item 再看 \texttt{rag/rag\_service.py}，理解检索结果如何交给模型总结。
\item 最后看 \texttt{prompts/*.txt} 和 \texttt{config/*.yml}，理解配置和提示词如何影响效果。
\end{enumerate}

\section{简历写法参考}

可以写成：

\begin{quote}
基于 LangChain + ReAct Agent + Chroma 构建企业知识库智能助手，支持企业制度问答、员工模拟数据查询和自动周报生成。系统通过 RAG 检索内部制度文档，结合 ReAct Agent 动态调用员工信息、报销记录、请假余额、项目进展等工具，并使用 Streamlit 实现流式对话、检索来源展示、知识文档上传和会话导出。
\end{quote}

项目亮点：

\begin{itemize}
\item 构建企业知识库 RAG 流程，支持 TXT/PDF 文档切分、向量化入库、top-k 检索和基于上下文的回答生成。
\item 基于 ReAct Agent 注册多类业务工具，实现制度问答、员工信息查询、项目记录查询和周报生成等多任务自动调度。
\item 使用 YAML 管理模型、Prompt、向量库和业务数据路径，提高项目可配置性和可迁移性。
\item 通过模拟企业数据集实现“知识库 + 结构化数据”的混合问答，增强项目业务完整度。
\item 构建 Streamlit 对话界面，支持模型选择、Temperature、Top-K、文件上传入库、引用来源展示和会话导出。
\end{itemize}

\section{验收问题与答案}

\subsection{制度类问题什么时候调用 \texttt{rag\_summarize}}

当用户询问 HR、财务报销、IT 支持、项目协作、信息安全、入职指引等公共制度或流程问题时，Agent 应优先调用 \texttt{rag\_summarize}，从企业知识库中获取依据后再回答。

\subsection{个人数据类问题为什么要先获取员工 ID}

因为年假余额、报销状态、项目记录、入职待办等数据都和具体员工绑定。用户问“我”的状态时，Agent 需要先调用 \texttt{get\_employee\_id}，再用 \texttt{fetch\_employee\_records} 查询模拟业务数据。

\subsection{周报生成需要哪些工具}

周报生成需要依次调用 \texttt{get\_employee\_id}、\texttt{get\_employee\_department}、\texttt{generate\_weekly\_report\_context}、\texttt{fetch\_employee\_records}，最后根据 \texttt{weekly\_report} 数据输出中文周报。

\subsection{如何让新增制度文档生效}

把 TXT/PDF 放入 \texttt{data/enterprise/} 后执行 \texttt{python -m rag.vector\_store}，或在 Streamlit 侧边栏上传文档并点击保存入库。入库成功后，新的制度内容会进入 Chroma 检索范围。

\section{一键运行流程}

首次运行：

\begin{Verbatim}[breaklines=true,breakanywhere=true,fontsize=\small]
cd D:\my_develop\A_work_program\LangChain-ReAct-Agent-main
pip install -r requirements.txt
$env:DASHSCOPE_API_KEY="your-api-key"
python -m rag.vector_store
streamlit run app.py
\end{Verbatim}

后续运行：

\begin{Verbatim}[breaklines=true,breakanywhere=true,fontsize=\small]
cd D:\my_develop\A_work_program\LangChain-ReAct-Agent-main
$env:DASHSCOPE_API_KEY="your-api-key"
streamlit run app.py
\end{Verbatim}

\section{注意事项}

\begin{itemize}
\item 当前员工数据是模拟数据，不接入真实公司系统。
\item \texttt{get\_employee\_id} 当前随机返回员工 ID，适合演示，不适合生产身份识别。
\item 项目用于学习、演示和简历展示时，不要放入真实员工隐私、客户信息、合同、源代码或其他敏感数据。
\item 如需完全重建向量库，请手动确认要处理的文件，不要使用批量递归删除命令。
\item \texttt{HELP.md} 当前不在项目中，本帮助文档直接维护 \texttt{HELP\_PROJECT\_GUIDE.tex}，避免旧生成脚本重新写入过期内容。
\end{itemize}

\section{参考资料}

\begin{itemize}
\item LangChain ReAct Agent API：\texttt{https://reference.langchain.com/python/langchain-classic/agents/react/agent/create\_react\_agent}
\item LangChain AgentExecutor API：\texttt{https://reference.langchain.com/python/langchain-classic/agents/agent/AgentExecutor}
\item LangChain Tools 文档：\texttt{https://docs.langchain.com/oss/python/langchain-tools}
\item LangChain Chroma 集成文档：\texttt{https://docs.langchain.com/oss/python/integrations/vectorstores/chroma/}
\item Chroma 官方文档：\texttt{https://docs.trychroma.com/docs/overview/getting-started}
\item Streamlit Chat Elements 文档：\texttt{https://docs.streamlit.io/library/api-reference/chat}
\item LangChain ChatTongyi 文档：\texttt{https://docs.langchain.com/oss/python/integrations/chat/tongyi}
\item LangChain DashScope Embeddings 文档：\texttt{https://docs.langchain.com/oss/python/integrations/text\_embedding/dashscope/}
\end{itemize}

\end{document}
